\chapter{Context, Requirements, and Project Framing}
\label{chap:context}

\section{Electricity consumption as a forecasting problem}

Electricity consumption is a time-dependent process whose distribution changes with the hour of day, day of week, season, household behavior, tariff incentives, and external conditions. Forecasting differs from ordinary regression because the ordering of observations is part of the problem. A model evaluated on a random sample of the same historical period may perform well while failing on the next day. The operational question is therefore conditional and time-indexed: given the measurements available up to time $t$, estimate consumption for times $t+1$ through $t+H$.

For a univariate household series $y_t$, a direct multi-horizon forecaster can be written as
\begin{equation}
  (\hat y_{t+1},\ldots,\hat y_{t+H}) = f_{\theta}(y_{t-L+1:t}, c_{t-L+1:t+H}),
  \label{eq:forecast}
\end{equation}
where $L$ is the lookback length, $H$ the forecast horizon, and $c$ represents calendar or other covariates known at forecast time. The packaged hourly TFTs use 336 hours and output 24 or 168 hourly values. The month model integrates persisted intervals into at least 270 valid rolling daily totals, uses up to 365 daily values, and outputs 30 daily values. All active models expose the 0.1, 0.5, and 0.9 quantiles.

This formulation must not be confused with a binary label such as ``high'' or ``low'' consumption. Classification can support anomaly or state detection, but it does not answer how many kWh will be consumed at each future hour. It can also become trivial if the class is constructed from a measurement retained among the inputs. One of the first framing decisions of this PFE was therefore to replace the historical class objective with a physically interpretable continuous target.

\section{Operational horizons and user value}

EnergyAI exposes three supported capabilities. The 24-hour horizon is intended for next-day review, and the 168-hour horizon supports week-ahead planning; both preserve hourly structure. The 30-day capability emits 30 daily totals rather than pretending that an integer duration of 720 means 720 hourly targets. It supports month-ahead budget planning and carries explicit daily cadence, target-count, input-readiness, and provenance metadata. Earlier monthly N-BEATS and MultiCycleNet results remain historical; the active month capability is a separately selected and audited Chronos-2 release.

The value of a forecast depends on its comparison target. Repeating yesterday's load or the same hour from the previous week often yields a surprisingly strong baseline. A deep model is only justified when it consistently improves over such rules, not merely when it reports a low number on an arbitrary normalized scale. This report consequently compares Global TFT with a seasonal-naive baseline and reports both average error and the proportion of households for which the model wins.

\section{Stakeholders}

The principal stakeholders are summarized in \cref{tab:stakeholders}. The platform is designed around an authenticated user who owns a site and its meter data. The administrator manages accounts and inspects system and model readiness. The supervisor and academic jury are evidence reviewers rather than product users: they need to verify the scientific narrative, result provenance, and limitations. Future energy managers or household users would require additional validation on their own local data before operational reliance.

\begin{table}[H]
\centering
\caption{Stakeholder needs and associated evidence.}
\label{tab:stakeholders}
\begin{tabularx}{\textwidth}{p{0.20\textwidth}YY}
\toprule
Stakeholder & Primary need & Evidence expected \\
\midrule
Normal user & Import consumption, inspect usage, generate supported forecasts, review actions & Clear data validation, forecast provenance, uncertainty, history, ownership isolation \\
Administrator & Operate accounts and inspect platform/model readiness & Role checks, health endpoints, audit events, artifact status \\
Supervisor and jury & Assess scientific rigor and engineering contribution & Reproducible protocol, primary citations, manifests, tests, Git timeline \\
Future operator & Deploy and monitor a stable system & Configuration, backups, dependency policy, observability, deployment validation \\
Data subject & Retain control over personal consumption data & Authentication, export, deletion, minimal exposure, bounded ingestion \\
\bottomrule
\end{tabularx}
\end{table}

\section{Functional scope}

The current repository supports the following primary user flow:
\begin{description}[style=multiline,leftmargin=1.8cm,labelwidth=1.4cm]
  \item[\textbf{FR-01}] create or access an account through the available authentication capability;
  \item[\textbf{FR-02}] complete site, timezone, tariff, and meter setup;
  \item[\textbf{FR-03}] preview and import a CSV containing timestamps and energy values;
  \item[\textbf{FR-04}] inspect consumption summaries and paginated readings;
  \item[\textbf{FR-05}] verify model readiness, choose 24 hours, 168 hours, or 30 daily targets, and generate a forecast with quantiles, cadence, and provenance;
  \item[\textbf{FR-06}] review alerts and recommendation-based actions;
  \item[\textbf{FR-07}] manage site, budget, account export, and account deletion settings; and
  \item[\textbf{FR-08}] permit administrators to manage users and inspect system/model readiness without exposing that surface to normal users.
\end{description}

Routes and endpoints exist for several secondary or gated behaviors, including Google authentication, email verification, password recovery, simulation controls, and administrative monitoring. This report describes them only to the degree supported by current code and configuration. It does not present removed navigation items or hidden research surfaces as primary Product V1 features.

\begin{figure}[H]
  \centering
  \includegraphics[width=0.98\textwidth]{use_case.pdf}
  \caption{EnergyAI use-case diagram derived from current frontend routes and backend routers.}
  \label{fig:usecase}
\end{figure}

\section{Non-functional requirements}

\subsection{Reliability and truthfulness}

The application must refuse unsupported operations rather than fabricate a result. Forecast readiness is a first-class contract. The backend checks artifact presence, checksum, input coverage, maximum consecutive gaps, finite values, and supported horizon before inference. Failure is returned as an explicit state that the frontend can explain.

\subsection{Privacy and isolation}

Electricity traces can reveal occupancy and routine. Each Site and Meter is associated with an owning user. API queries must apply ownership conditions, and administrative access must be explicit. Authentication uses short-lived access tokens and refresh credentials stored as hashes server-side; the browser receives refresh state through an HttpOnly cookie rather than a JSON token saved in web storage.

\subsection{Latency and resource constraints}

The user experience requires bounded imports and responsive inference. The product therefore packages fixed CPU-loadable artifacts instead of retraining in response to user requests. The two TFT checkpoints are approximately 5.6 MB each. The month LoRA adapter is approximately 4.9 MB, but its pinned 120M-parameter base adds a substantial image-cache and memory cost. A measured CPU cold load took about 9.0 seconds; warm 30-day inference took about 77 ms. The production image therefore caches the exact base revision and verifies both base and adapter checksums before offline serving.

\subsection{Cold-start generalization}

Known-household performance can hide dependence on household identity. A product that accepts a new user's data must also be assessed on households that contributed no gradient updates. The Low Carbon London experiments split entire households into known and cold-start cohorts before model training. Checkpoint-selection claims in this report prioritize cold-start macro metrics, even when a global metric appears higher.

\subsection{Reproducibility}

Every central model claim must map to a manifest or leaderboard, and every packaged checkpoint must map to a SHA-256 digest. The report itself maintains registries for models, figures, notebooks, commits, artifact hashes, tests, and requirement-to-test traceability. This requirement is stricter than merely retaining source code: it connects the exact file and experimental status to each published claim.

\section{System boundary and assumptions}

EnergyAI retains a reproducible six-service Docker Compose configuration for local development and controlled demonstrations. The jury release was updated through Git commit \texttt{6796b1ec7360f9c05cb804fbd8be29ba1ea78ef8} and deployed to Microsoft Azure on 9 August 2026. The verified cloud boundary comprises public Next.js and FastAPI Container Apps, a dedicated email-worker Container App, Azure Database for PostgreSQL, Azure Container Registry, and Azure Key Vault. Azure-managed hostnames provide HTTPS for the web and API endpoints. Public email registration and verification, SMTP delivery, and Google identity are enabled and smoke-tested. The deployed worker scope remains narrower than local Compose: alert and avatar-cleanup workers are absent, and avatar storage is ephemeral. Likewise, validation metrics describe London and benchmark cohorts; public availability does not guarantee performance on a Moroccan household or organization without a local evaluation.

The cover page uses the official combined UMI/FSM logo and blue-orange motif published by the Faculty of Sciences of Meknes. The department and programme wording were verified against the current institutional programme pages, and the faculty identifies the academic supervisor as Pr. Ali Oubelkacem. No official defense-date evidence is preserved in the project, so the date remains an explicit administrative placeholder. The publicly available cover model is intended for doctoral theses rather than Master PFE reports; consequently, the English PFE wording is an evidence-based adaptation, while the jury composition also remains pending confirmation.

\section{Project risks and acceptance criteria}

\begin{table}[H]
\centering
\caption{Principal risks and acceptance criteria.}
\label{tab:risks}
\begin{tabularx}{\textwidth}{p{0.24\textwidth}YY}
\toprule
Risk & Control & Acceptance evidence \\
\midrule
Target leakage & Causal feature audit and chronological split & Reproducible comparison between historical and corrected pipelines \\
Misleading model comparison & Common cohorts, units, horizons, and baselines & Leaderboards with explicit protocol columns \\
Unsupported product promise & Capability, cadence, and readiness gates & 24 and 168 hourly targets plus exactly 30 daily targets \\
Cross-user data exposure & Ownership-scoped queries and RBAC & PostgreSQL integration and authorization tests \\
Corrupted or mismatched model & Manifest and SHA-256 verification & Strict load, finite output, and expected shape checks \\
Research--serving transform mismatch & Explicit normalization contracts and frozen-cohort serving re-evaluation & Transform comparison completed; new-site efficacy remains conditional \\
Dependency or browser gate failure & Automated quality and security audits & Current QA report including conditional findings \\
Institutional non-compliance & Official template and administrative review & Placeholder closure before final submission \\
\bottomrule
\end{tabularx}
\end{table}

This framing makes the scientific and product objectives mutually reinforcing. Reliable software cannot repair an invalid experiment, while a valid model cannot create user value without data controls, understandable output, and operational evidence.
